\chapter{Two Roads Between Newton and the Action}
\label{two-roads}

\begin{quote}\itshape
Every textbook on classical mechanics must decide what is a postulate and
what is a theorem. The choice is rarely made explicit, yet it determines which
phenomena look natural and which look like paradoxes.
\end{quote}

\section{Introduction: the same theory, two logical orders}
\label{two-orders}

Classical mechanics is usually presented in one of two logical orders.

\begin{description}
  \item[Newton first.] The postulates are the second law
  $m\ddot{\mathbf r}_a=\mathbf F_a$ (together with a catalogue of force laws
  and a statement about constraint reactions). Lagrange's equations, the
  principle of least action and the Hamiltonian formalism are \emph{derived}.
  This is the order of Goldstein \cite{Goldstein}, Arnold's \cite{Arnold-cm} first chapters, and most
  introductory courses.
  \item[Action first.] The postulate is Hamilton's principle,
  $\delta S=0$ with $S=\int L\,dt$. Newton's equations are \emph{derived}, and
  conservation laws follow from symmetries of $L$ via Noether's theorem. This
  is the order of Landau and Lifshitz \cite{Landau-Lifshitz-1}, 
        Lanczos \cite{Lanczos} and of much of modern field theory.
\end{description}

Both orders lead to the same equations for simple systems of point particles, 
so it is
tempting to regard the choice as a matter of taste. It is not. The two roads
are \emph{not} mutually inverse once we leave the simplest class of systems:
the road from Newton to the action requires hypotheses (potential forces,
ideal constraints) that the road back silently assumes, and the two roads give
\emph{different} answers for nonholonomic constraints. This chapter walks both
roads carefully, marks the places where hypotheses enter, and compares them
as pedagogical strategies.

In the language of chapter \ref{rm}, the Newtonian laws are
statements anchored in the real world: forces are measured, constraints are
realized by hinges, rails and blades. The action is a construction of the ideal
world: a functional whose extremals are, by design, trajectories. The
paradoxes of this chapter all occur at the interface where one is converted
into the other.

\section{From Newton to the action}
\label{newton-to-action}

\subsection{d'Alembert's principle: the bridge}

Consider $N$ particles with positions $\mathbf r_a$, masses $m_a$, subject to
applied forces $\mathbf F_a$ and constraint reactions $\mathbf R_a$:
\begin{equation}
  m_a\ddot{\mathbf r}_a=\mathbf F_a+\mathbf R_a .
  \label{newton}
\end{equation}
The reactions $\mathbf R_a$ are unknown. Newtonian mechanics closes the system
only by supplementing~\eqref{newton} with a statement about the
\emph{nature} of constraint forces. The statement that makes the passage to
Lagrangian mechanics possible is the following.

\begin{quote}
\textbf{Ideal constraints.} The reactions do no work on any \emph{virtual
displacement}, that is, on any displacement $\delta\mathbf r_a$ compatible with
the constraints at a frozen instant of time:
\begin{equation}
  \sum_a \mathbf R_a\cdot\delta\mathbf r_a = 0 .
  \label{ideal}
\end{equation}
\end{quote}

Substituting~\eqref{ideal} into~\eqref{newton} eliminates the unknown
reactions and gives d'Alembert's principle:
\begin{equation}
  \sum_a\bigl(m_a\ddot{\mathbf r}_a-\mathbf F_a\bigr)\cdot\delta\mathbf r_a=0
  \qquad\text{for all virtual displacements.}
  \label{dalembert}
\end{equation}
It is worth stressing that~\eqref{ideal} is \emph{not} a theorem of
Newtonian mechanics. It is a constitutive assumption about hinges, smooth
surfaces and rolling without slipping. For a frictionless rail it is true by
the definition of ``frictionless''; for rolling contact it is an empirical
statement about the contact force doing no work at the instantaneous point of
contact; for a rough surface it is false.

\subsection{Generalized coordinates and Lagrange's equations}

If the constraints are \emph{holonomic}, $f_k(\mathbf r_1,\dots,\mathbf r_N,t)=0$,
we may choose independent generalized coordinates $q^1,\dots,q^n$ with
$\mathbf r_a=\mathbf r_a(q,t)$. Virtual displacements are then arbitrary
$\delta q^i$, and~\eqref{dalembert} becomes
\begin{equation}
  \sum_i\Bigl[\frac{d}{dt}\frac{\partial T}{\partial\dot q^i}
  -\frac{\partial T}{\partial q^i}-Q_i\Bigr]\delta q^i=0,
  \qquad Q_i=\sum_a\mathbf F_a\cdot\frac{\partial\mathbf r_a}{\partial q^i},
  \label{lagrange-T}
\end{equation}
the key algebraic step being the identity
$\frac{d}{dt}\frac{\partial\mathbf r_a}{\partial q^i}
=\frac{\partial\dot{\mathbf r}_a}{\partial q^i}$. Independence of the
$\delta q^i$ gives $n$ equations. If moreover $Q_i=-\partial V/\partial q^i$
(or, more generally, $Q_i=\frac{d}{dt}\frac{\partial U}{\partial\dot q^i}
-\frac{\partial U}{\partial q^i}$ for a generalized potential $U$), then with
$L=T-V$ we obtain
\begin{equation}
  \frac{d}{dt}\frac{\partial L}{\partial\dot q^i}-\frac{\partial L}{\partial q^i}=0 .
  \label{EL}
\end{equation}

\subsection{From Lagrange's equations to an action}

Equation~\eqref{dalembert} is a statement at one instant. To turn it into
a statement about whole trajectories, integrate over $[t_0,t_1]$ and integrate
the term $m_a\ddot{\mathbf r}_a\cdot\delta\mathbf r_a$ by parts, using
$\delta\mathbf r_a(t_0)=\delta\mathbf r_a(t_1)=0$:
\begin{equation}
  \int_{t_0}^{t_1}\Bigl(\delta T+\sum_a\mathbf F_a\cdot\delta\mathbf r_a\Bigr)dt=0 .
  \label{integral-dalembert}
\end{equation}
This is the \emph{Lagrange--d'Alembert principle}. It is valid for \emph{any}
applied forces, including friction and driving. Only when
$\sum_a\mathbf F_a\cdot\delta\mathbf r_a=-\delta V$ (or $\delta U$ with the
appropriate velocity-dependent form) does the integrand become an exact
variation and~\eqref{integral-dalembert} collapse into Hamilton's principle
\begin{equation}
  \delta S=0,\qquad S[q]=\int_{t_0}^{t_1}L(q,\dot q,t)\,dt .
  \label{hamilton}
\end{equation}

This is the first moral of the Newton-first road. The \emph{variational}
principle is the special case of a more general \emph{virtual-work}
principle, valid when the active forces are derivable from a (generalized)
potential. Dissipative forces fit into~\eqref{integral-dalembert} but not
into~\eqref{hamilton}; this is why friction is a nuisance for the
action-first road and a non-event for the Newton-first road.

\subsection{Constraints on the Newton-first road}

On this road constraints cost nothing conceptually. There are two equivalent
techniques.

\begin{enumerate}
  \item \textbf{Elimination.} Solve the constraints, pass to $q^i$,
  and write~\eqref{EL}. The reactions never appear.
  \item \textbf{Multipliers.} Keep redundant coordinates $x^\mu$ and impose
  $f_k(x)=0$. Because virtual displacements satisfy
  $\partial_\mu f_k\,\delta x^\mu=0$, the vector $\mathbf R$ in~\eqref{ideal}
  must be a combination of the gradients $\partial_\mu f_k$:
  \begin{equation}
    m\ddot x_\mu=F_\mu+\sum_k\lambda_k\,\partial_\mu f_k .
    \label{multipliers}
  \end{equation}
  The multipliers $\lambda_k$ are \emph{the constraint reactions}, with direct
  physical meaning (tension in a rod, normal force on a surface).
\end{enumerate}

The multiplier technique has a pedagogical advantage that is usually
squandered: the multipliers answer questions about the real world
(will the rope go slack? will the bead leave the surface?) that the reduced
Lagrangian cannot.

\begin{example}[Plane pendulum]\label{pendulum}
Take coordinates $(x,y)$ with $y$ upward and the constraint
$f=x^2+y^2-\ell^2=0$. Equation~\eqref{multipliers} reads
\[
  m\ddot x=2\lambda x,\qquad m\ddot y=-mg+2\lambda y .
\]
Writing $x=\ell\sin\theta$, $y=-\ell\cos\theta$ and projecting on the radial
and tangential directions gives
\[
  \ell\ddot\theta=-g\sin\theta,\qquad
  T:=-2\lambda\ell=m\bigl(\ell\dot\theta^{2}+g\cos\theta\bigr).
\]
The first is what elimination gives directly. The second is the tension, and
the criterion $T>0$ tells us when a pendulum on a string (rather than a rod)
stops behaving like a pendulum.
\end{example}

\section{From the action back to Newton}
\label{action-to-newton}

Now reverse the order. Postulate~\eqref{hamilton}, with $L=T-V$ and
Cartesian coordinates for a free particle in a potential.

\subsection{The derivation, and what it uses}

Let $q(t)\to q(t)+\varepsilon\,\eta(t)$ with $\eta(t_0)=\eta(t_1)=0$. Then
\[
  \frac{dS}{d\varepsilon}\Big|_0=\int_{t_0}^{t_1}
  \Bigl(\frac{\partial L}{\partial q}-\frac{d}{dt}\frac{\partial L}{\partial\dot q}\Bigr)\eta\,dt=0
  \quad\text{for all }\eta .
\]
The \emph{fundamental lemma of the calculus of variations} (the integrand is
continuous and orthogonal to every test function, so it vanishes) gives the
Euler--Lagrange equation~\eqref{EL}. For $L=\tfrac12m\dot q^2-V(q)$ this is
$m\ddot q=-V'(q)$: Newton's second law.

Four assumptions have been used silently.
\begin{enumerate}
  \item The functional \emph{has} the form $T-V$. Nothing in the principle
  says so. Nature (or the teacher) supplies $L$.
  \item Variations are taken with fixed endpoints in \emph{time and
  position}. The motion is therefore determined by a boundary-value problem,
  although Newton's law is an initial-value problem.
  \item The word ``least'' is not justified. The principle is
  \emph{stationarity}.
  \item The variations $\eta$ range over all smooth functions: there are no
  constraints on the admissible paths. Constraints are discussed in
  Section~\ref{constraints-variational}.
\end{enumerate}

\subsection{Stationary, not least}
\label{least}

The traditional name ``principle of least action'' is misleading. For the
harmonic oscillator $\ddot q=-\omega^2q$ the action of the true trajectory
between $t_0$ and $t_1$ is a minimum if and only if $t_1-t_0<\pi/\omega$;
beyond that the second variation acquires a negative direction, and the
true path is a saddle. In general, a trajectory ceases to be a local minimum
at the first \emph{conjugate point} (Jacobi). The geometric optics analogue is
a caustic. On a sphere, the great circle between two points is minimal only
up to the antipode.

This matters pedagogically. A student told that ``the particle chooses the
path of least action'' will naturally ask how it \emph{knows} the final point
in advance. The honest answer is that it does not: the boundary-value
formulation is an equivalent \emph{global description} of an \emph{initial-value}
motion, and stationarity is the local property that survives. This is the
point at which the action-first road most invites teleological language,
precisely the sort of mixing of real-world causality and ideal-world logic
that the  chapter \ref{rm} of this book warn against.

\subsection{The inverse problem and non-uniqueness}
\label{inverse}

The road from the equations of motion to a Lagrangian is not a function. Three
observations:
\begin{itemize}
  \item $L$ and $L+\frac{d}{dt}F(q,t)$ give identical equations (the action
  changes by boundary terms only).
  \item $L$ and $cL$ give identical equations for any constant $c\neq0$.
  \item Genuinely different Lagrangians can give the same equations. The
  two-dimensional isotropic oscillator, for example, admits Lagrangians that
  are not related by a total derivative or a constant factor.
\end{itemize}
The existence question (``is a given system of second-order equations
Euler--Lagrange?'') is answered by the Helmholtz conditions: in particular a
system $\ddot q^i=f^i(q,\dot q,t)$ with friction linear in the velocity,
$\ddot q=-\gamma\dot q-\omega^2q$, is Euler--Lagrange only after an explicit
time-dependent factor $e^{\gamma t/m}$ is introduced (the Caldirola--Kanai
Lagrangian), and then the canonical momentum no longer equals the mechanical
momentum. The dissipative world is not forbidden to the action, but it is not
natural to it.

\section{Constraints on the action-first road}
\label{constraints-variational}

\subsection{Holonomic constraints: agreement}

Suppose we vary $S$ over paths that satisfy $f_k(x)=0$ at every time.
Introduce the augmented action
\begin{equation}
  \tilde S=\int\bigl[L(x,\dot x)+\lambda_k(t)\,f_k(x)\bigr]dt .
  \label{aug-holonomic}
\end{equation}
Its Euler--Lagrange equations are precisely the multiplier
equations~\eqref{multipliers}, with $\lambda_k$ as the (time-dependent)
reaction forces. For holonomic constraints the two roads agree completely, and
the multipliers play the same physical role.

A conceptual remark: in the action-first road the multipliers are
\emph{introduced to enforce a restriction on trial paths}. In the Newton-first
road they are the \emph{forces that actually do the enforcing}. That the two
coincide is a theorem, not a tautology; it follows from~\eqref{ideal}.

\subsection{Nonholonomic constraints: disagreement}
\label{nonholonomic}

A nonholonomic constraint restricts velocities, not positions:
\begin{equation}
  a^k_i(q)\,\dot q^i=0,\qquad k=1,\dots,m,
  \label{nonhol}
\end{equation}
where the one-forms $a^k=a^k_i\,dq^i$ do not integrate to the form $df_k$.
Two natural generalizations of the holonomic procedure exist, and they give
different equations (see section \ref{cm:nh}) .

\paragraph{Route A: Lagrange--d'Alembert (Newton-first).}
Take~\eqref{ideal} as the definition of an ideal constraint and impose the
constraint on virtual displacements:
$a^k_i\,\delta q^i=0$. Then
$\bigl(\frac{d}{dt}\frac{\partial L}{\partial\dot q^i}-\frac{\partial L}{\partial q^i}\bigr)\delta q^i=0$
for all admissible $\delta q$, so
\begin{equation}
  \frac{d}{dt}\frac{\partial L}{\partial\dot q^i}-\frac{\partial L}{\partial q^i}
  =\lambda_k\,a^k_i .
  \label{LdA}
\end{equation}
This is \emph{not} the Euler--Lagrange equation of any augmented action built
from~\eqref{nonhol}; it is a virtual-work principle that does not
integrate into~\eqref{hamilton}.

\paragraph{Route B: vakonomic (action-first).}
Treat the problem as a variational problem with a nonholonomic side
condition. Extremize $S$ among paths satisfying~\eqref{nonhol}; the
multipliers enter as in~\eqref{aug-holonomic},
$\tilde L=L+\lambda_k a^k_i\dot q^i$, and the Euler--Lagrange equations
for $q^i$ are
\begin{equation}
  \frac{d}{dt}\frac{\partial L}{\partial\dot q^i}-\frac{\partial L}{\partial q^i}
  =-\dot\lambda_k\,a^k_i-\lambda_k\bigl(\partial_ja^k_i-\partial_ia^k_j\bigr)\dot q^j .
  \label{vako}
\end{equation}

Comparing~\eqref{LdA} and~\eqref{vako}: the term $-\dot\lambda_ka^k_i$
alone could be absorbed by redefining the multiplier, which is what happens
in the holonomic case. The second term involves the \emph{curvature} of the
constraint distribution, $d a^k=\tfrac12(\partial_ja^k_i-\partial_ia^k_j)\,dq^j\wedge dq^i$,
and is nonzero precisely when the constraint is not integrable. In other
words:
\begin{quote}
\emph{The two procedures coincide if and only if the constraint is
holonomic.}
\end{quote}
The disagreement is not a computational subtlety. It is a statement about
how nature behaves, and nature has voted against vakonomic (section \ref{cm:nh}).

\begin{example}[A skate]\label{skate}
A skate is a rigid body on a plane whose blade touches the ice at its center
of mass. Coordinates $(x,y,\varphi)$, with Lagrangian
$L=\tfrac12m(\dot x^2+\dot y^2)+\tfrac12I\dot\varphi^2$, and the blade cannot
slip sideways:
\[
  a=\sin\varphi\,dx-\cos\varphi\,dy,\qquad a_\varphi=0 .
\]
\emph{Route~A.} The reaction is perpendicular to the blade and applied at the
center of mass, so it exerts no torque: $I\ddot\varphi=0$ and the speed along
the blade is constant. This is what a skater coasting on a straight line
observes.

\emph{Route~B.} The $\varphi$-component of~\eqref{vako} reads
\[
  I\ddot\varphi=\lambda\,v,\qquad v=\dot x\cos\varphi+\dot y\sin\varphi,
\]
so in general the ``constraint'' exerts a torque about the center of mass
which the real blade cannot exert. Moreover, Route~B is a
\emph{boundary-value} problem: for given endpoints the multiplier is fixed
by the final configuration, and the particle's initial velocity can no longer be
prescribed freely.
\end{example}

Experiments on rolling bodies, skates and sleighs (and the engineering
practice of wheeled vehicles) are consistent with Route~A. Route~B is not
wrong, but it answers a different question: it describes the optimal control
of a system whose velocities are \emph{steered} (a sub-Riemannian
geodesic problem, with applications to robotics and parking a car), not the
free motion of a system whose velocities are \emph{forced} by a passive
constraint. The principle of least action is therefore a theorem about
\emph{potential} dynamics under \emph{holonomic} constraints, and a
definition or a hypothesis beyond that.

\subsection{Constraints as limits}
\label{constraints-limits}

A third answer to ``what is a constraint?'' replaces the idealization by a
family of smooth models. A holonomic constraint $f=0$ is the limit
$\kappa\to\infty$ of a stiff potential $V_\kappa=\tfrac12\kappa f^2$. The limit
is delicate (see chapter \ref{lm}). If the initial position lies on the surface
and the initial
velocity is tangent to it, the motion converges to the d'Alembert motion on the
surface (Rubin and Ungar). If the initial velocity has a normal component, the
normal motion is a fast oscillation with a conserved action, and the limiting
motion on the surface is governed by the original Hamiltonian \emph{plus an
extra potential} that depends on that action. Rubin and Ungar 
\cite{RubinUngar1957}proved this for
codimension one; Takens \cite{Takens1980} treated the general case, 
under a non-resonance
condition on the frequencies of the normal oscillations. When that condition
fails (frequencies crossing along the surface), there may be no differential
equation on the surface governing the limiting motion at all.
Similarly, a nonholonomic constraint of the type~\eqref{nonhol} is
the limit of a large viscous friction $-\mu\,(a_i\dot q^i)\,a^i$ as
$\mu\to\infty$, and \emph{this} limit reproduces Route~A (the
Lagrange--d'Alembert equations), not Route~B.

This is a clean illustration of the thesis that a theory is not a limit of
another but a chart on it: the same rod, with the same symbol $f=0$, belongs
to different physical charts depending on whether it is the limit of a stiff
spring or the limit of a viscous friction.

\section{Comparing the routes as pedagogy}
\label{pedagogy}

\begin{table}[ht]
\centering\small
\renewcommand{\arraystretch}{1.25}
\begin{tabular}{@{}p{3.1cm}p{5.6cm}p{5.6cm}@{}}
\toprule
 & \textbf{Newton $\to$ action} & \textbf{Action $\to$ Newton}\\
\midrule
Postulates & $m\ddot{\mathbf r}=\mathbf F$, ideal constraints, force laws &
  $\delta S=0$, form of $L$, admissible paths \\
Primary objects & forces, reactions, accelerations & energies, symmetries, trajectories as a whole \\
Constraint forces & computed directly; physical meaning of $\lambda_k$ & introduced as multipliers; meaning proved later \\
Dissipation & trivial (virtual-work form) & awkward (needs doubled variables, time-dependent $L$, or is excluded) \\
Nonholonomic constraints & natural (Route~A) & ambiguous (Route~B differs) \\
Conservation laws & derived case by case & Noether's theorem: unified \\
Generalization to fields and QM & indirect & direct (path integral, field actions) \\
Typical misconception & ``the Lagrangian is just Newton in disguise'' & ``the particle chooses the least-action path'' \\
\bottomrule
\end{tabular}
\caption{Logical structure of the two routes.}
\label{routes}
\end{table}

\subsection{Newton first}

\emph{Strengths.} Every step has a physical referent. Students see where
the action comes from and, crucially, \emph{when it does not exist}.
Constraint reactions are computed, not suppressed. The nonholonomic
and dissipative cases are included without a change of philosophy.

\emph{Weaknesses.} The passage to generalized coordinates is algebraically
heavy, and the key identity
$\frac{d}{dt}\partial_{q}\mathbf r=\partial_{\dot q}\dot{\mathbf r}$ looks
like a trick. The reader emerges thinking Lagrangian mechanics is a
\emph{computational device}, and misses that the action is a primary
object in field theory and quantum mechanics.

\subsection{Action first}

\emph{Strengths.} Symmetry becomes the organizing principle:
Noether's theorem explains energy, momentum and angular momentum at once.
Coordinate independence is manifest: $\delta S=0$ does not know what
coordinates are. The road to field theory and quantization is short.

\emph{Weaknesses.} The phrase ``least action'' invites teleology
(Section~\ref{least}). The student learns that dissipative systems are
``not mechanics'' and that constraint reactions are an afterthought. Most
damagingly, the nonholonomic case produces a \emph{wrong answer} that is
indistinguishable, in form, from the right one; a student who has only
been taught the variational principle will confidently apply the vakonomic
equations to a rolling ball.

\subsection{Conclusions}

Both routes coincide for restricted but very large and important class of
holonomic systems.  Newton-first route is probably better when classical 
mechanics is seen as a part of CNT. Since this book is devoted to second
role of classical mechanics as a base for quantum and relativistic theory
we'll take the action-first route.
